% GENERATED FILE. Edit canonical lessons, then run npm run generate:books.
% canonical-source-hash: fnv1a64:62e70ce918e0bfa6
% canonical-lessons: SA-C43-hospitality, SA-C43-conduct, SA-C43-humility, SA-C43-trust, SA-C43-honest

\chapter{Good Manners}
\label{ch:sa-good-manners}
\hlchaptermodality{43}

\begin{chapteropening}
\textbf{By the end of this chapter:} \emph{I can name hospitality, good conduct, humility, trust and plain dealing in Sanskrit.}

Complete SA-C43-honest and demonstrate the chapter promise.
\end{chapteropening}

\section[ātithyam]{\sk{आतिथ्यम्} (ātithyam) — hospitality}

\label{lesson:SA-C43-hospitality}



\begin{quote}
Before the new word: what did \emph{bahu} mean?
\end{quote}

\subsection*{You'll want to know: \sk{आतिथ्यम्}}
\textbf{\sk{आतिथ्यम्}} (\emph{ātithyam}) — \textquotedblleft{}hospitality\textquotedblright{}.

\textbf{\sk{आतिथ्यम्}} is hospitality: not the guest, but everything owed to one.

The \sk{अतिथिः} you already know is stretched into the treatment that guest receives, so the word names a duty rather than a person.

Classically it runs in an order, and you have the words for all of it: an \sk{आसनम्} first, then \sk{पानीयम्}, then \sk{अन्नम्}. \textbf{\sk{आतिथ्यम्}} is that whole sequence in one word.

\begin{grammarlens}[title={What you've built}]
The guest word, turned into what is owed to the guest.
\end{grammarlens}

\subsection*{Your turn}
Say these aloud:

\begin{itemize}
\raggedright
  \item \emph{ātithyam}
  \item it once more, slowly
  \item \emph{bahu}, then \emph{ātithyam}
\end{itemize}

\begin{tcolorbox}[breakable,colback=teal!4,colframe=teal!35!black,title={Before you move on}]
What does \sk{आतिथ्यम्} mean? (\textquotedblleft{}hospitality\textquotedblright{}.) Say it once more.
\end{tcolorbox}
\section[vinayaḥ]{\sk{विनयः} (vinayaḥ) — good conduct, training}

\label{lesson:SA-C43-conduct}



\begin{quote}
Before the new word: what did \emph{ātithyam} mean?
\end{quote}

\subsection*{You'll want to know: \sk{विनयः}}
\textbf{\sk{विनयः}} (\emph{vinayaḥ}) — \textquotedblleft{}good conduct, training\textquotedblright{}.

\textbf{\sk{विनयः}} is the bearing of someone well brought up — modest, disciplined, easy to be near.

Its pieces are \emph{vi-}, \textquotedblleft{}away, apart\textquotedblright{}, and the leading doing-word \sk{नयति}: to have \textbf{\sk{विनयः}} is to have been led away from what was rough in you. Training and good manners are one word here, because the second is thought to come out of the first.

The Buddhist books of monastic training carry this word as their name.

\begin{grammarlens}[title={What you've built}]
Manners understood as training rather than as polish.
\end{grammarlens}

\subsection*{Your turn}
Say these aloud:

\begin{itemize}
\raggedright
  \item \emph{vinayaḥ}
  \item it once more, slowly
  \item \emph{ātithyam}, then \emph{vinayaḥ}
\end{itemize}

\begin{tcolorbox}[breakable,colback=teal!4,colframe=teal!35!black,title={Before you move on}]
What does \sk{विनयः} mean? (\textquotedblleft{}good conduct, training\textquotedblright{}.) Say it once more.
\end{tcolorbox}
\section[namratā]{\sk{नम्रता} (namratā) — humility}

\label{lesson:SA-C43-humility}



\begin{quote}
Before the new word: what did \emph{vinayaḥ} mean?
\end{quote}

\subsection*{You'll want to know: \sk{नम्रता}}
\textbf{\sk{नम्रता}} (\emph{namratā}) — \textquotedblleft{}humility\textquotedblright{}.

\textbf{\sk{नम्रता}} is humility — bendingness, the quality of one who bows.

It stands on \emph{nam}, \textquotedblleft{}to bow\textquotedblright{}, which is the \emph{nam} that opens \sk{नमस्ते} on your first page and sits inside \sk{प्रणामः}. Those two are single acts; \textbf{\sk{नम्रता}} is the same bow turned into a settled habit.

Hindi \emph{namratā} is the word people still use of someone who carries success lightly.

\begin{grammarlens}[title={What you've built}]
The bow you already make, become a habit.
\end{grammarlens}

\subsection*{Your turn}
Say these aloud:

\begin{itemize}
\raggedright
  \item \emph{namratā}
  \item it once more, slowly
  \item \emph{vinayaḥ}, then \emph{namratā}
\end{itemize}

\begin{tcolorbox}[breakable,colback=teal!4,colframe=teal!35!black,title={Before you move on}]
What does \sk{नम्रता} mean? (\textquotedblleft{}humility\textquotedblright{}.) Say it once more.
\end{tcolorbox}
\section[viśvāsaḥ]{\sk{विश्वासः} (viśvāsaḥ) — trust, confidence}

\label{lesson:SA-C43-trust}



\begin{quote}
Before the new word: what did \emph{namratā} mean?
\end{quote}

\subsection*{You'll want to know: \sk{विश्वासः}}
\textbf{\sk{विश्वासः}} (\emph{viśvāsaḥ}) — \textquotedblleft{}trust, confidence\textquotedblright{}.

\textbf{\sk{विश्वासः}} is trust — and the picture inside it is breathing out.

It is \emph{vi-} with \emph{śvas}, \textquotedblleft{}to breathe\textquotedblright{}: to trust someone is to let the breath go, to stop holding it. English says much the same thing when it calls a friend a confidant, from Latin for \emph{with faith}.

Hindi \emph{viśvās} is unchanged, and \emph{viśvāsghāt}, a betrayal, is the same word with the breaking of it attached.

\begin{grammarlens}[title={What you've built}]
Trust, described as the moment you stop holding your breath.
\end{grammarlens}

\subsection*{Your turn}
Say these aloud:

\begin{itemize}
\raggedright
  \item \emph{viśvāsaḥ}
  \item it once more, slowly
  \item \emph{namratā}, then \emph{viśvāsaḥ}
\end{itemize}

\begin{tcolorbox}[breakable,colback=teal!4,colframe=teal!35!black,title={Before you move on}]
What does \sk{विश्वासः} mean? (\textquotedblleft{}trust, confidence\textquotedblright{}.) Say it once more.
\end{tcolorbox}
\section[saralaḥ]{\sk{सरलः} (saralaḥ) — straightforward, honest}

\label{lesson:SA-C43-honest}



\begin{quote}
Before the last word of the chapter: what did \emph{namratā} mean, and what did \emph{viśvāsaḥ} mean?
\end{quote}

\subsection*{You'll want to know: \sk{सरलः}}
\textbf{\sk{सरलः}} (\emph{saralaḥ}) — \textquotedblleft{}straightforward, honest\textquotedblright{}.

\textbf{\sk{सरलः}} is straight: of a \sk{मार्गः}, of a piece of \sk{दारु}, and of a person who says the thing plainly.

The pine is called \emph{sarala} because it grows straight up, and the same word does duty for a road with no bends and for someone with no side to them.

English runs the metaphor the same way — \emph{straight}, \emph{upright}, \emph{a straight answer} — which is worth noticing, because the two languages arrived at it separately.

\begin{grammarlens}[title={What you've built}]
Five words for how a person carries themselves: hospitality, conduct, humility, trust, plain dealing.
\end{grammarlens}

\subsection*{Your turn}
Say these aloud:

\begin{itemize}
\raggedright
  \item \emph{saralaḥ}
  \item it once more, slowly
  \item all five in order, then \emph{ātithyam} and \emph{saralaḥ} together
\end{itemize}

\begin{tcolorbox}[breakable,colback=teal!4,colframe=teal!35!black,title={Before you move on}]
What does \sk{सरलः} mean? (\textquotedblleft{}straightforward, honest\textquotedblright{}.) Say it once more.
\end{tcolorbox}
